\section{Pair closure, a positive-vector witness, and renewal}
\label{s:recursion}

All networks in this section are finite plane networks with the site
laws of Definition~\ref{d:moment} and the original legs of
Definition~\ref{d:pinmodel}.  The boundary used in every distance is
the original protected set \(S=S(G)\).  Ordinary domains have the
meaning of Definition~\ref{d:domain}.  We first close the response
class under tracing, and then derive a nonnegative shift recursion.
The quantitative hypotheses are (H1) and (H2) of
Definitions~\ref{h:H1} and~\ref{h:H2}, and the positive-vector
assertion (H3\(_h\)) defined below.

The trace operations, the stopped trail class, and the coefficients
\(a_E(w),R_E^\bullet(w)\) are those of
Definition~\ref{d:trace}.  The ratios of
Definition~\ref{d:ratios} telescope to
\(R_E^\bullet(w)=\prod_j\rho_j^\bullet\), with
\(\rho_j^\tau=\rho_j^0=1\) at every admissible
\(\mathrm D\) departure.  Proposition~\ref{p:pairwalk} supplies
the finite expansion, its target constants, and its separate-degree
bound.  We use it only with its prescribed stopping set
\(S\cup\{\text{target}\}\) and check its start and target
hypotheses at both expansions below.

The local comparison used here is Lemma~\ref{l:local}, including
its uniformity in outside spins, original unit pins, missing
external factors, and actual tip data allowed by the ordinary-domain
rules.  The ordered product estimate is Lemma~\ref{l:exposure};
the common pointwise majorant and its summed consequence are
Lemma~\ref{l:onepoint}.  These are results of the preceding walk
section.  In particular the last lemma bounds the sum of pointwise
maxima for the two pin choices and permits its twelve exceptional
departures at arbitrary trail positions.

\begin{definition}\label{d:pairenv}
For \(D\in\mathcal D_{\rm cons}(G)\), let \(P_B(D)\) and \(P_C(D)\)
be the ordered nonremoved pairs \((u,v)\) at distance two,
collinear or diagonal respectively, with a designated common neighbor
\(c\) of status \(\mathrm X\). Each endpoint is \(\mathrm U\),
or \(\mathrm D\) with at most one present bond.
For integer \(d\ge0\), let \(F_X(d)\) be the supremum of
\[
 \left|\mathbb E_D^\tau(n_u\cdot n_v)
                    -\mathbb E_D^0(n_u\cdot n_v)\right|
\]
over every finite \(G\), every original unit-pin assignment,
every \(D\in\mathcal D_{\rm cons}(G)\), and every pair in \(P_X(D)\)
with \(\min\{d(u,S(G)),d(v,S(G))\}\ge d\).
Use supremum zero for an empty class, and set \(F_X(d)=F_X(0)\)
for negative integers. Thus \(0\le F_X\le2\), and \(F_X\)
is nonincreasing.
\end{definition}

Indeed, both expectations are expectations under positive finite
probability laws by Lemma~\ref{l:finitefacts}, and
\(|n_u\cdot n_v|\le1\).  Increasing \(d\) restricts the class over
which the supremum is taken.  In particular the asserted bounds and
monotonicity also hold when that class is empty.

\begin{lemma}\label{l:closure}
Let \(D\in\mathcal D_{\rm cons}(G)\), \((u,v)\in P_X(D)\) have
center \(c\), and \(u\notin S\). Before tracing, if \(u\) is
\(\mathrm U\), drop its kept bond to the removed \(c\); this changes
no partition function. Trace any trail from \(u\) without visiting
\(S\) before its last site; it need not avoid \(v\).
After every prefix ending at the live endpoint \(y\):
all departed sites are \(\mathrm D\) or \(\mathrm X\);
consistency between nonremoved sites holds; all nonlive
\(\mathrm U\) sites have full kept sets and all nonlive \(\mathrm D\)
sites have exactly two kept bonds.
After a positive-length prefix a live \(\mathrm U\) site has lost
only its arrival bond, and a live \(\mathrm D\) site keeps one bond;
the arrival site's other endpoint is \(\mathrm D\) or \(\mathrm X\).
At length zero a live \(\mathrm U\) site has its three non-\(c\)
kept bonds, while a live \(\mathrm D\) site retains its two kept
bonds but has at most one present bond.
Whenever \(y\notin S\), removing \(y\) yields an ordinary member
of \(\mathcal D_{\rm cons}(G)\). Every \(\mathrm U\) departure
at \(y\) before stopping has a covered pattern. In the ordinary
domain obtained by removing the live site \(y\) from the pre-step
network, every ordered pair of distinct nonremoved non-arrival
neighbors of \(y\) belongs to \(P_B\) or \(P_C\), with designated
center \(y\).
If the endpoint is \(e\in S\), it is \(\mathrm U\) with just its
arrival bond lost. This is the sole kept-set defect in the terminal
auxiliary network; all original legs remain. No removal of \(e\)
or ordinary-domain membership after removing \(e\) is asserted.
\end{lemma}

\begin{proof}
Write the trail as \((v_0,\ldots,v_m)\), with \(v_0=u\).
All departures, arrivals, and removals in the following induction
are the operations of Definition~\ref{d:trace}.
Every neighbor of a site outside \(S\) is in \(G\), since such a
site has distance at least three from \(\mathbb Z^2\setminus G\).
Thus the three non-\(c\) kept bonds mentioned at length zero really
exist as bonds inside \(G\).  A bond incident to a removed site is
absent regardless of its kept flag at the other end.  Dropping the
flag at \(u\) for \(\{u,c\}\) therefore changes no factor, pinned
or free, and changes no normalized expectation.  The initial live
site is now the only possible failure of the ordinary kept-set
conditions.  Consistency is unchanged because \(c\) is removed.

We prove the assertions by induction on the prefix length.  At a
departure from a live \(\mathrm U\) site, the incoming bond is absent
and its three other bonds are kept.  The outgoing bond is present and
differs from the incoming bond.  Dropping the outgoing bond and
changing the law to \(\lambda_{\mathrm D}\) leaves exactly two kept
bonds at the departed site.  If the initial site is \(\mathrm D\),
it has at most one present bond, which must be the chosen outgoing
bond.  Its removal therefore leaves no other present bond incident
to it.  At any later live \(\mathrm D\) site there is just one kept
bond, and an actual departure uses that bond and removes the site.
These are all possible departure cases.  In particular a departed
site never becomes \(\mathrm U\) again.

On arrival the new live site cannot be removed, since the traversed
bond was present.  If it is \(\mathrm U\), it was not previously
departed from; before this arrival its kept set was full.  Dropping
the arrival bond leaves precisely the stated live kept set.  If it
is \(\mathrm D\), its nonlive kept set had two elements and contained
the present arrival bond; dropping that bond leaves one element.
All traversed bonds are dropped by both ends.  Removing a
\(\mathrm D\) departure site destroys only bonds to a removed site,
and so cannot create a consistency failure between nonremoved sites.
Every other kept flag is unchanged.  This proves the induction,
including the assertion about the other end of the arrival bond.

Before stopping, no protected site has been visited.  Removing the
live endpoint when it is outside \(S\) removes the only exceptional
kept set; all protected sites still have their original full kept
sets.  The remaining \(\mathrm U\) kept sets are full, the remaining
\(\mathrm D\) kept sets have size two, and consistency holds.
Consequently the resulting network is an ordinary member of
\(\mathcal D_{\rm cons}(G)\).  At a terminal hit \(e\in S\),
the site \(e\) has not previously been visited.  It still has its
original \(\mathrm U\) law and has dropped just the arrival bond.
All other sites satisfy the ordinary kept-set conditions.  Tracing
never changes an original leg factor, including those at \(e\).

It remains to verify the local pattern and pair assertions.  At a
\(\mathrm U\) departure \(y\), the designated arrival site is
\(c\) at the initial departure, and is the preceding trail site
otherwise.  It is removed, or is \(\mathrm D\) with its two kept
bonds among the three bonds other than its bond to \(y\).
Every nonremoved non-arrival neighbor of \(y\) keeps its bond to
\(y\): the live \(\mathrm U\) site keeps that bond, so consistency
forces the other end to keep it.  Such a neighbor of status
\(\mathrm D\) has exactly one additional kept bond.  Every corner
of status \(\mathrm D\) has two kept bonds, and every other
\(\mathrm U\) site has its full kept set.  The outgoing neighbor
is nonremoved because its bond is present.  These are exactly the
non-tip rules of Definition~\ref{d:pattern}.  The diamond of radius
two about \(y\) lies in \(G\).  Original legs incident to its
corners or tips occupy external slots.  Actual missing external
factors, both slots at each corner, and actual tip data are treated
by the covering procedure of Lemma~\ref{l:local}; that procedure
does not require consistency with an actual tip kept set.

Finally let \(p\) be a pair of distinct nonremoved non-arrival
neighbors of \(y\).  The two sites are at distance two and are
collinear or diagonal.  Removing \(y\) supplies their designated
removed center.  Each endpoint of status \(\mathrm D\) previously
had its bond to \(y\) present, so after removal it has at most one
present bond.  The endpoints therefore meet
Definition~\ref{d:pairenv} in the ordinary domain just obtained.
None of this reasoning used avoidance of \(v\).
\end{proof}


\begin{lemma}\label{l:terminalcover}
In the first expansion of a pair of Definition~\ref{d:pairenv},
assume also that \(v\notin S\), and
suppose the trail avoids its target \(v\) and stops at \(e\in S\).
In the terminal auxiliary network, start a second expansion from \(v\)
to \(S\), again dropping its kept bond to \(c\) if it is \(\mathrm U\).
The starts of both expansions have their designated incoming bond
absent and are admissible in Definition~\ref{d:trace}.
Before the second trail hits \(S\), a \(\mathrm U\) departure
outside \(\{y':0<d(y',e)\le2\}\) has a covered pattern.
There are at most twelve exceptional \(\mathrm U\) departures
on each second trail. Both expansions satisfy Lemma~\ref{l:onepoint}.
\end{lemma}

\begin{proof}
The first expansion has an admissible start outside \(S\) by
Proposition~\ref{p:pairwalk}.  Its incoming bond is \(\{c,u\}\),
which is absent because \(c\) is removed.  At \(u\), a
\(\mathrm U\) law has at most three present bonds; a
\(\mathrm D\) law has at most one by the pair definition.

By Lemma~\ref{l:closure}, the terminal network has precisely the
fixed defect at \(e\) in addition to the ordinary data elsewhere.
The first trail never visits \(v\), so its status and kept set
have not changed.  Its number of present bonds can only have
decreased.  Since \(v\notin S\), it is an admissible second start:
it is \(\mathrm U\) with at most three present bonds, or
\(\mathrm D\) with at most one.  Its designated incoming bond is
\(\{c,v\}\), again absent.  Dropping that kept flag at a
\(\mathrm U\) start leaves every partition function unchanged.
Thus neither expansion can initially reverse across its designated
incoming bond.

For completeness the second pair expansion, with target \(e\),
also meets the target degree hypothesis in
Proposition~\ref{p:pairwalk}.  Before its first arrival, the
protected site \(e\) had at most four factors in total, including
original legs.  The first arrival dropped one present bond, leaving
at most three factors.  Its law remains \(\lambda_{\mathrm U}\).
The second start is distinct from \(e\) because \(v\notin S\).

Repeat the kept-set induction of Lemma~\ref{l:closure} during the
second trace, retaining \(e\) as a fixed defective site.  Before a
hit of \(S\), no second departure is at \(e\), and no other
protected site changes its kept set.  Outside \(e\) and the live
endpoint, the full \(\mathrm U\) and two-element \(\mathrm D\)
kept-set rules continue to hold.  Consistency holds for bonds not
incident to \(e\).  This follows in each departure and arrival
case from the same two flag changes as in the preceding induction;
no operation changes a leg factor.

If \(d(y',e)>2\), the defective site is outside the template
about a \(\mathrm U\) departure \(y'\).  Its non-tip data
therefore satisfy every rule in Definition~\ref{d:pattern}.
Any absent external factor is allowed by Lemma~\ref{l:local},
including a factor whose absence is caused by the defect outside
the template.  The actual tips are also distinct from \(e\),
so their statuses and kept sets satisfy that lemma's ordinary-domain
tip rules.  Every outside spin has norm at most one; the original
leg factors are retained.  No ordinary-domain assumption is imposed
on the whole terminal auxiliary network.  Since no departure is
at \(e\), every possible
exception lies in
\(\{y':0<d(y',e)\le2\}\).  There are four sites at distance
one and eight at distance two in the square lattice.  A site can
be departed from with status \(\mathrm U\) at most once, since
that departure changes its status to \(\mathrm D\).  Hence there
are at most twelve exceptions, at arbitrary positions along the
second trail.

The first trace has no exceptions by Lemma~\ref{l:closure},
even if its stopping rule is enlarged to all first hits of \(S\).
Both traces thus have the absent incoming bond, admissible start,
and exception bound required by Lemma~\ref{l:onepoint}.  Whenever
(H1) and (H2) are assumed, its estimate applies to each trace.
\end{proof}

\begin{remark}
The additional hypothesis \(v\notin S\) in
Lemma~\ref{l:terminalcover} is necessary.  For example, take
\(G=\{-3,\ldots,3\}^2\), so
\(S=\{x\in G:\|x\|_\infty\ge2\}\).  Give \(c=(1,0)\)
status \(\mathrm X\) and every other site status \(\mathrm U\)
with its full kept set.  The pair \(u=(0,0),v=(2,0)\) belongs to
\(P_B\), and the first trail
\(((0,0),(0,1),(0,2))\) avoids \(v\) and stops in \(S\).
Nevertheless \(v\in S\), so it cannot be the start of a second
trace of Definition~\ref{d:trace}.  In the recursion below both
endpoints have distance at least one from \(S\), as required.
\end{remark}

\begin{lemma}\label{l:log-mean-bounds}
Let \(\mathcal L(x,y)=(x-y)/(\log x-\log y)\) for \(x\ne y\)
and \(\mathcal L(x,x)=x\), for positive \(x,y\).
Put \(f(s)=1/\mathcal L(s,1)\), so \(f(1)=1\).
For \(3/5\le c\le1\), the function
\[
 G_c(x,y)=\frac1{3\mathcal L(cx,y)}-\frac1{5\mathcal L(x,y)}
\]
is nonincreasing in each positive variable. In particular,
for \(x\ge a=116/125\), \(y\ge1\),
\[
 G_c(x,y)\le f(ac)/3-f(a)/5.
\]
For rational \(0<s<1\), put \(t=(1-s)/(1+s)\) and
\[
 f_-(s)=\frac2{1-s}\sum_{k=0}^{7}\frac{t^{2k+1}}{2k+1},
 \qquad
 f_+(s)=f_-(s)+\frac{2t^{17}}{17(1-t^2)(1-s)}.
\]
Then \(f_-(s)\le f(s)\le f_+(s)\).
For the two rows of the following table, all numbers are exact rationals:
\[
 \begin{array}{c|c|ccc}
 \tau&c_\tau&a_\tau&b_\tau&b'_\tau\\ \hline
 {\rm s}&961/1000&145217/10^6&352782/10^6&340008/10^6\\
 {\rm t}&187/200&150000/10^6&357549/10^6&344661/10^6
 \end{array}
\]
They satisfy
\[
 \frac{f_+(ac_\tau)}3-\frac{f_-(a)}5\le a_\tau,\qquad
 \frac{f_+(ac_\tau)}3\le b_\tau,\qquad
 \frac{f_+(c_\tau)}3\le b'_\tau,\qquad
 0<b'_\tau\le b_\tau .
\]
\end{lemma}

\begin{proof}
For positive \(r,s\), direct integration, also when \(r=s\), gives
\begin{equation}\label{e:rec-log-integral}
 \frac1{\mathcal L(r,s)}
       =\int_0^1\frac{\mathrm dt}{(1-t)r+ts}.
\end{equation}
In particular the reciprocal logarithmic mean is decreasing in each
argument, and \(\mathcal L\) is increasing in each argument.
For fixed \(t\), put \(A=(1-t)x\) and \(B=ty\).
The integrand for \(G_c\) is
\(1/(3(cA+B))-1/(5(A+B))\).  Its derivatives with respect to
\(x\) and \(y\) are respectively
\[
 (1-t)\left(\frac1{5(A+B)^2}
                   -\frac{c}{3(cA+B)^2}\right),\qquad
 t\left(\frac1{5(A+B)^2}
                   -\frac1{3(cA+B)^2}\right).
\]
Both are nonpositive.  Indeed \(cA+B\le A+B\), \(5c\ge3\),
and hence
\[
 5c(A+B)^2\ge3(cA+B)^2,\qquad
 5(A+B)^2\ge3(cA+B)^2.
\]
On every compact positive rectangle the denominators are bounded
away from zero, justifying differentiation under the integral in
\eqref{e:rec-log-integral}.  This proves both monotonicities and
the value at the lower endpoints \((a,1)\).

For the logarithm enclosure, \(0<t<1\) and
\(-\log s=\log((1+t)/(1-t))\).  Integrating the identity
\[
 \frac1{1-z^2}=\sum_{k=0}^7z^{2k}+\frac{z^{16}}{1-z^2}
       \qquad(0\le z\le t)
\]
from zero to \(t\) shows that
\[
 0\le-\log s-2\sum_{k=0}^7\frac{t^{2k+1}}{2k+1}
   =2\int_0^t\frac{z^{16}}{1-z^2}\,\mathrm dz
   \le\frac{2t^{17}}{17(1-t^2)}.
\]
Division by \(1-s>0\) proves the asserted enclosure, with no
approximate logarithms.

Here are explicit rational checks for all the numerical endpoints.
Let \(K=10^9\).  The integers in the last two columns satisfy
\(\ell_s/K\le f_-(s)\le f_+(s)\le u_s/K\):
\begin{equation}\label{e:rec-log-table}
\begin{array}{c|c|r|r}
 s&t&\ell_s&u_s\\ \hline
 116/125&9/241&1037827030&1037827031\\
 27869/31250&3381/59119&1058344574&1058344575\\
 5423/6250&827/11673&1072644315&1072644316\\
 961/1000&39/1961&1020022307&1020022308\\
 187/200&13/387&1033980764&1033980765
\end{array}
\end{equation}
For clarity, the table is checked solely by the following integer
formula.  If \(s=u/v\), \(t=p/q\) are the fractions in a row, put
\[
 \begin{gathered}
 D_*=45045=\operatorname{lcm}(1,3,5,7,9,11,13,15),\\
 P(p,q)=\sum_{k=0}^7\frac{D_*}{2k+1}p^{2k+1}q^{14-2k},
 \qquad Q=(v-u)D_*q^{15}.
 \end{gathered}
\]
All these denominators are positive integers, and substitution gives
\[
 f_-(s)=\frac{2vP(p,q)}{Q},\qquad
 f_+(s)=\frac{2v\{17(q^2-p^2)P(p,q)+D_*p^{17}\}}
                    {17(q^2-p^2)Q}.
\]
Thus the two checks in each row are the explicit integer inequalities
\[
 \ell_sQ\le2vKP(p,q),\qquad
 2vK\{17(q^2-p^2)P(p,q)+D_*p^{17}\}
                \le17u_s(q^2-p^2)Q.
\]
Inserting the five displayed pairs \((u/v,p/q)\) verifies each
inequality by sums and products of nonnegative integers.
To check the requested coefficient bounds, write
\(A_\tau=10^6a_\tau\), \(B_\tau=10^6b_\tau\), and
\(B'_\tau=10^6b'_\tau\).  After using
\eqref{e:rec-log-table}, their positive residual numerators are
\[
\begin{array}{c|r|r|r}
 \tau&15000A_\tau-5u_{ac_\tau}+3\ell_a
       &3000B_\tau-u_{ac_\tau}&3000B'_\tau-u_{c_\tau}\\ \hline
 {\rm s}&13215&1425&1692\\
 {\rm t}&259510&2684&2235
\end{array}
\]
The respective positive denominators are \(15K,3K,3K\).
The two values of \(b'_\tau\) are positive and at most the
corresponding \(b_\tau\) by direct comparison of their integer
numerators.  This proves every table inequality.
\end{proof}

\begin{lemma}\label{l:response-typed}
Fix \(c_{\rm s}=961/1000\), \(c_{\rm t}=187/200\) and assume (H1).
At a covered \(\mathrm U\) departure of type \(\tau\in\{{\rm s},{\rm t}\}\),
the response coefficients in Lemma~\ref{l:response} satisfy
\[
 0<\alpha_j\le a_\tau,\qquad 0<\beta_j\le b_\tau
       \quad\hbox{with three present neighbors},\qquad
 0<\beta'_j\le b'_\tau\quad\hbox{with two}.
\]
Here the exact constants are those of Lemma~\ref{l:log-mean-bounds};
in particular \(a_{\rm t}=150000/10^6\).
With \(a=116/125\), their analytic bounds are
\[
 \beta_j\le f(ac_\tau)/3,\qquad
 \beta'_j\le f(c_\tau)/3,\qquad
 \alpha_j\le f(ac_\tau)/3-f(a)/5.
\]
These follow using \(I^\tau\ge c_\tau J^\tau\),
\(I^0\ge J^0\ge1\), and \(J^\tau\ge a\).
Also \(a_\tau\le18094/100000\) and
\(b_\tau,b'_\tau\le36751/100000\).
\end{lemma}

\begin{proof}
Use \(E'\), \(I,J\) from Definition~\ref{d:ratios}, and abbreviate
\(I^\bullet=\mathbb E_{E'}^\bullet I\),
\(J^\bullet=\mathbb E_{E'}^\bullet J\).  A covered departure is
outside \(S\), as in Lemma~\ref{l:local}.  That lemma and (H1)
give \(I^\tau/J^\tau\ge c_\tau\); its stated uniformity covers
the actual tips and all missing external factors, including both
corner slots.  We derive the coefficient identities directly from
Definition~\ref{d:ratios}.  For positive \(A,B\), including
\(A=B\),
\(\log A-\log B=(A-B)/\mathcal L(A,B)\).  Since
\(\rho_j^\bullet=J^\bullet/I^\bullet\), with three present
neighbors this gives
\[
 \begin{split}
 \log\rho_j^0-\log\rho_j^\tau
 &=\frac{I^\tau-I^0}{\mathcal L(I^\tau,I^0)}
       -\frac{J^\tau-J^0}{\mathcal L(J^\tau,J^0)}\\
 &=\left(\frac1{3\mathcal L(I^\tau,I^0)}
            -\frac1{5\mathcal L(J^\tau,J^0)}\right)r_{fg}
       +\frac{r_{fo}+r_{go}}{3\mathcal L(I^\tau,I^0)}.
 \end{split}
\]
Thus the response of Lemma~\ref{l:response} has the canonical
coefficients
\begin{equation}\label{e:rec-typed-coefficients}
 \beta_j=\frac1{3\mathcal L(I^\tau,I^0)},\qquad
 \alpha_j=\frac1{3\mathcal L(I^\tau,I^0)}
                     -\frac1{5\mathcal L(J^\tau,J^0)}.
\end{equation}
This fixes the coefficient choice used below, also when responses
vanish or logarithmic-mean arguments coincide.

With three present neighbors, free positivity gives
\[
 J^0\ge1,\qquad
 I^0-J^0=\tfrac2{15}\mathbb E^0_{E'}P_{fg}
             +\tfrac13\mathbb E^0_{E'}(P_{fo}+P_{go})\ge0.
\]
These uses of Lemma~\ref{l:finitefacts} satisfy its degree
condition: after \(y\) is removed, each neighbor has at most
three incident factors and no original leg, so insertion of one
scalar product has separate degree at most four.  The sites
\(f,g,o\) are pairwise nonadjacent.  On every positive
conditioning fiber of the other spins, their spins are independent
and their conditional means have norm at most \(3/5\) by
Lemmas~\ref{l:finitefacts} and~\ref{l:35}.  Hence
\(\mathbb E^\tau_{E'}P_{fg}\ge-9/25\), and
\[
 J^\tau=1+\tfrac15\mathbb E^\tau_{E'}P_{fg}
          \ge1-\tfrac9{125}=a.
\]
The same conditioning under either pin choice yields
\[
 \frac{41}{50}\le I^\bullet\le\frac{34}{25},\qquad
 \frac{116}{125}\le J^\bullet\le\frac{134}{125}.
\]
Indeed, for the three conditional means \(m_f,m_g,m_o\),
\[
 -\frac{27}{50}\le
 \sum_{\{b,b'\}\subset\{f,g,o\}}m_b\cdot m_{b'}
 =\frac{|m_f+m_g+m_o|^2-|m_f|^2-|m_g|^2-|m_o|^2}{2}
 \le\frac{27}{25}.
\]
Together with \(|m_f\cdot m_g|\le9/25\), averaging proves
both intervals.  Formula~\eqref{e:rec-log-integral} places
\(\mathcal L\) between its positive arguments, so the very same
canonical coefficients satisfy
\[
 \alpha_j\ge\frac{25}{102}-\frac{25}{116}
              =\frac{175}{5916}>0,\qquad \beta_j>0.
\]
All four arguments of the logarithmic means are strictly positive.
By their monotonicity and \(I^\tau\ge c_\tau J^\tau\),
\(I^0\ge J^0\), formula~\eqref{e:rec-typed-coefficients} yields
\[
 \beta_j\le\frac1{3\mathcal L(c_\tau a,1)},\qquad
 \alpha_j\le G_{c_\tau}(J^\tau,J^0)
        \le\frac{f(ac_\tau)}3-\frac{f(a)}5.
\]
The last step uses \(3/5\le c_\tau\le1\) and
Lemma~\ref{l:log-mean-bounds}.

With two present neighbors, \(J=1\) and
\(I=1+P_{fo}/3\).  Free positivity gives \(I^0\ge1\),
while local (H1) gives \(I^\tau\ge c_\tau\).
The same divided-difference calculation now gives the response
\(\beta'_jr_{fo}\), where
\(\beta'_j=1/(3\mathcal L(I^\tau,I^0))>0\), so
\(\beta'_j\le f(c_\tau)/3\).  With just one present neighbor
both ratios are one by Definition~\ref{d:ratios}, and there is
no response coefficient.  At a \(\mathrm D\) departure that
definition likewise makes both ratios one, so its response is zero.
The rational enclosures of Lemma~\ref{l:log-mean-bounds} prove
the stated bounds.  Finally the largest \(a_\tau\) is
\(150000/10^6<18094/100000\), and each \(b_\tau,b'_\tau\)
is at most \(357549/10^6<36751/100000\).
\end{proof}

\begin{definition}\label{d:kernel}
For \(X=B,C\), fix \(c=0\), \(u=(-1,0)\), and \(v=(1,0)\)
for \(B\), \(v=(0,1)\) for \(C\).
Let \(\mathcal T_n^X\) be the length-\(n\) plane trails avoiding \(c\),
from \(u\) to their first hit of \(v\), with incoming direction \(u-c\).
A departure is a lattice revisit when its site was departed from earlier.
At a lattice first visit \(i\), put \(x_i=x_{\tau_i}\) according to
its straight/turn type. Of the three neighbors other than the arrival
site, write \(o\) for the outgoing neighbor and \(f,g\) for the others.
A neighbor is dead if it is \(c\) or was departed from twice before \(i\).
The outgoing neighbor of a true trail is not dead.
If \(f,g\) are both live, assign coefficient \(a_{\tau_i}\) to
\(\{f,g\}\), and \(b_{\tau_i}\) to each of \(\{f,o\},\{g,o\}\).
If exactly one of \(f,g\) is live, assign \(b'_{\tau_i}\) to its
pair with \(o\); if neither is live there is no response pair.
The constants are those of Lemma~\ref{l:log-mean-bounds}.
A live pair \(p\) has type \(Y=B,C\) and shift
\[
 s(p)=\max_{z\in p}\min\{d(z,u),d(z,v)\}.
\]
Let \(\kappa_{Y,s}(i)\) sum its assigned coefficients.
For any fixed, \(\theta\)-independent orderings \(\sigma_w\) of lattice
first visits, define
\[
 \begin{split}
 M_{XY,s}
 &=\sum_{n\ge1}\sum_{w\in\mathcal T_n^X}
     3^{-n}(3/5)^{\operatorname{rev}(w)}
     \sum_k\kappa_{Y,s}(\sigma_w(k))
                   \prod_{l\le k}x_{\sigma_w(l)},\\
 M_\theta[X,Y]&=\sum_{s\ge0}M_{XY,s}\theta^{-s}.
 \end{split}
\]
These are nonnegative series. The identity
\(3^{-n}(3/5)^r=3^{-(n-r)}5^{-r}\) retains the mass-free
trace coefficients. This \(M\) is the type-dependent response kernel.
\end{definition}

Until finiteness is proved, each series is understood in
\([0,+\infty]\) as the supremum of its finite subsums.
Regrouping these nonnegative sums therefore requires no prior
real-valued summability assertion.

To justify the assertion about the outgoing neighbor, a site other
than \(u\) that has already been departed from twice has used three
distinct incident bonds: two outgoing bonds and an intervening
arrival bond.  Entering it again uses its fourth bond, leaving no
possible subsequent departure.  It cannot be the final target,
since the target has no earlier departure.  The start \(u\) has
only three available bonds because \(c\) is excluded; after two
departures all three have been used.  Neither case permits an
outgoing step to a dead site on a trail in \(\mathcal T_n^X\).
The incoming direction also forbids the initial reversal, since that
step would enter \(c\).

\begin{definition}\label{d:kernel-order}
Fix
\[
 h=(1,h_C),\quad h_C=\frac{5404319552844595}{4503599627370496},
 \qquad \rho_0=\frac{993015}{10^6},\qquad L=12 .
\]
For a first visit set \(\kappa_h(i)=\sum_{Y,s}\kappa_{Y,s}(i)h_Y\).
For each \(X\in\{B,C\}\) and \(w\in\mathcal T_n^X\), define
\(\sigma_w\) as follows. If \(n\le L\), order its lattice
first-departure indices by increasing
\((x_i-1)/(x_i\kappa_h(i))\), assigning value \(+\infty\) when
\(\kappa_h(i)=0\), and break ties by increasing path index.
If \(n>L\), list straight first departures in path order and
then turning first departures in path order.
Hereafter \(M_{XY,s}\) and \(M_\theta\) use exactly these orderings.
They do not depend on \(\theta\).
\end{definition}

Both \(x_{\rm s}\) and \(x_{\rm t}\) exceed one, and \(h_C>1\).
Thus every finite ordering key has a positive denominator.  The
infinite-key rule and the index tie-break give a unique ordering on
each finite trail.  All data in these comparisons are rational.

\begin{definition}[(H3\(_h\))]\label{h:H3}
For the type-dependent kernel and fixed orderings of
Definitions~\ref{d:kernel}--\ref{d:kernel-order}, (H3\(_h\)) is the
following mathematical assertion, with the specified \(h,\rho_0\):
for \(X=B,C\) and every integer \(R_0\ge0\),
\[
 \sum_{n=1}^{R_0}\sum_{w\in\mathcal T_n^X}
 3^{-n}(3/5)^{\operatorname{rev}(w)}
 \sum_k\kappa_h(\sigma_w(k))\prod_{l\le k}x_{\sigma_w(l)}
       \le\rho_0h_X .
\]
Equivalently, the full nonnegative sums satisfy \(M_1h\le\rho_0h\).
The assertion is proved using the stationary lattice Bellman
certificates, with no periodic spatial walk or remainder.
\end{definition}

The equivalence follows directly by taking increasing finite
length cutoffs.  For each length there are finitely many trails,
and each trail has finitely many first visits and response pairs.
Rearrangement by \((Y,s)\) preserves a nonnegative sum.  Consequently
the supremum of the displayed finite sums is exactly
\(\sum_{Y,s}M_{XY,s}h_Y\).  A bound for the full sum bounds each
cutoff, and bounds for every cutoff bound their supremum.

\begin{lemma}\label{l:Mtheta}
Assume (H2) and (H3\(_h\)) for the type-dependent kernel.
For every \(q^2<\theta\le1\), all entries of \(M_\theta\) are
finite. There exist \(q^2<\theta<1\) and \(\rho_0<\rho<1\)
such that the same vector \(h\) satisfies \(M_\theta h\le\rho h\).
Every shift of a length-\(n\) trail is at most \(n/2+1\), and
each matrix entry is bounded by
\[
 \sum_{n\ge1} n(\bar\alpha+2\bar\beta)
                      C_V q^n\theta^{-n/2-1}.
\]
\end{lemma}

\begin{proof}
Let \(w=(v_0,\ldots,v_n)\in\mathcal T_n^X\).  At departure
index \(i\),
\[
 d(v_i,u)\le i,\qquad d(v_i,v)\le n-i.
\]
The minimum of these two distances is at most \(n/2\).  Each
site in a response pair is a neighbor of \(v_i\); the triangle
inequality therefore gives \(s(p)\le n/2+1\).
Put \(C=\bar\alpha+2\bar\beta>0\).  At each first visit, the
sum of all its response coefficients is at most \(C\), since
the rational coefficient table has \(a_\tau\le\bar\alpha\) and
\(b_\tau,b'_\tau\le\bar\beta\).  There are at most \(n\) first
visits.  Since every \(x_i\ge1\), an initial product in any
ordering is at most the product over all first visits.  Hence
the length-\(n\) contribution to a fixed entry of \(M_\theta\)
is at most
\[
 nC\theta^{-n/2-1}
       \sum_{w\in\mathcal T_n^X}\mathsf W(w)
       \le nCC_Vq^n\theta^{-n/2-1}.
\]
The weights of Definition~\ref{d:weights} agree because
\(3^{-n}(3/5)^{\operatorname{rev}(w)}\prod_i x_i
 =\mathsf W(w)\).  These trails form a subset of the open trails
with the same initial data in (H2).  Since
\(q/\sqrt\theta<1\), the displayed series converges; explicitly
it equals
\(CC_V\theta^{-1}(q/\sqrt\theta)/(1-q/\sqrt\theta)^2\).
This part uses only (H2).

We now give an explicit choice transporting the witness.  Set
\[
 H=\frac{\max h}{\min h},\quad
 b=\frac{1+q}{2},\quad r=\frac qb,\quad
 B_*=2HCC_Vb^{-2}\frac{r}{(1-r)^3},
\]
and define
\begin{equation}\label{e:rec-shift-choice}
 \varepsilon=\min\left\{\frac{1-b}{2},
                         \frac{1-\rho_0}{2B_*}\right\},
 \qquad \theta=1-\varepsilon,\qquad
 \rho=\frac{1+\rho_0}{2}.
\end{equation}
Here \(0<q<b<1\), \(0<r<1\), \(B_*>0\), and
\(0<\varepsilon\le(1-b)/2\).  Thus
\(q^2<q<b<\theta<1\), and \(\rho_0<\rho<1\).
For every integer shift \(0\le s\le n/2+1\le n+1\),
the mean value theorem on \([\theta,1]\) gives
\[
 0\le\theta^{-s}-1
     \le(1-\theta)s b^{-s-1}
     \le(1-\theta)(n+1)b^{-n-2}.
\]
This includes \(s=0\), whose left side is zero.
For a fixed row \(X\), multiply each coefficient by
\(h_Y/h_X\le H\), use the same length bound, and sum.  All
summands are nonnegative and finite by the preceding estimate, so
\[
 \frac{((M_\theta-M_1)h)_X}{h_X}
 \le (1-\theta)HCC_Vb^{-2}
                  \sum_{n\ge1}n(n+1)r^n
 = (1-\theta)B_*\le\frac{1-\rho_0}{2}.
\]
The series identity follows by differentiating the geometric series
twice, or by multiplying its partial sums by \((1-r)^3\) and
taking the limit.  Adding \((M_1h)_X/h_X\le\rho_0\) proves
\(M_\theta h\le\rho h\) for the same \(h\) and orderings.
\end{proof}

\begin{lemma}\label{l:recineq}
Fix the two H1 thresholds above. Assume (H1) and (H2), use
the bounds of Lemma~\ref{l:response-typed}, and use the
type-dependent kernel with the fixed orderings of
Definition~\ref{d:kernel-order}.
For \(X=B,C\) and every integer \(d\ge1\),
\[
 F_X(d)\le\sum_{Y=B,C}\sum_{s\ge0}M_{XY,s}F_Y(d-s)+C_Tq^{2d},
 \qquad
 C_T=2(12/5)^{24}\left(\frac{C_V}{1-q}\right)^2 .
\]
\end{lemma}

\begin{proof}
Fix a finite \(G\), unit pins \(\tau\), an ordinary consistent
domain \(D\), and \((u,v)\in P_X(D)\) with designated center
\(c\), such that
\[
 d_0=\min\{d(u,S),d(v,S)\}\ge d\ge1.
\]
In particular \(u,v\notin S\).  Normalize the kept flag from
\(u\) to \(c\) when \(u\) is \(\mathrm U\), as in
Lemma~\ref{l:closure}.  We continue to denote this normalized
starting network by \(D\); all its partition functions and
expectations agree with those of the original ordinary domain.
Both endpoints meet the restrictions of Proposition~\ref{p:pairwalk}:
a \(\mathrm U\) endpoint has no original leg and at most three
present bonds because \(c\) is removed, and a \(\mathrm D\)
endpoint has at most one present bond by the pair definition.

Fix once and for all a plane lattice symmetry \(\varphi\) sending
the designated triple \((c,u,v)\) to the canonical triple for
type \(X\) in Definition~\ref{d:kernel}.  Explicitly, for \(B\)
choose a signed coordinate permutation \(Q\) with
\(Q(u-c)=(-1,0)\); then \(v-c=-(u-c)\) implies
\(Q(v-c)=(1,0)\).  For \(C\), the vectors \(u-c,v-c\) are
perpendicular lattice unit vectors, so choose the signed coordinate
permutation with \(Q(u-c)=(-1,0)\) and \(Q(v-c)=(0,1)\).
In both cases set \(\varphi(z)=Q(z-c)\).  This choice depends
only on the original ordered pair and its designated center,
and is used for every trail in both sums below.  In particular
\(w\mapsto\varphi(w)\) is injective.

Use that proposition for both choices \(\bullet=\tau,0\), with
stopping set \(S\cup\{v\}\).  The set of trails and their
coefficients \(a_D(w)\) depend only on the status and bond data,
and hence are common to the two expansions.  Split them into
\(\mathcal W_v\), ending at \(v\), and \(\mathcal W_S\),
ending at a point of \(S\) while avoiding \(v\).  The terminal
factor for \(w\in\mathcal W_v\) is a common
\(\eta_v\in\{1,3/5\}\), specifically \(1\) for a
\(\mathrm U\) target and \(3/5\) for a \(\mathrm D\) target,
by Proposition~\ref{p:pairwalk}.  For \(w\in\mathcal W_S\) ending
at \(e\), put
\(X^\bullet(w)=\mathbb E_{D[w]}^\bullet(n_e\cdot n_v)\).
The response equals \(T_1+T_2\), where
\begin{align*}
 T_1&=\sum_{w\in\mathcal W_v}a_D(w)
                 (R_D^\tau(w)-R_D^0(w))\eta_v,\\
 T_2&=\sum_{w\in\mathcal W_S}a_D(w)
              \{R_D^\tau(w)X^\tau(w)-R_D^0(w)X^0(w)\}.
\end{align*}
These sums are finite: every trail uses distinct bonds from a
fixed finite network.

For \(T_1\), fix a trail \(w=(v_0,\ldots,v_n)\in\mathcal W_v\).
Let \(\Phi_U(w)\) be its departure indices whose current status
is \(\mathrm U\).  At each such index Lemmas~\ref{l:closure}
and~\ref{l:local} give a covered pattern and
\(\rho_i^\tau\le x_{\tau_i}\).  Its lattice type is computed
using the same incoming direction as the trace, including
\(u-c\) at the initial departure.  Order \(\Phi_U(w)\) by
restricting \(\sigma_w\), after identifying the pair with the
canonical pair by the fixed symmetry \(\varphi\); thus here
\(\sigma_w\) means the ordering of \(\varphi(w)\).
The trail is admissibly started with absent incoming bond
\(\{c,u\}\), and \(x_{\tau_i}\ge1\).  These are precisely the
start and product-bound hypotheses of Lemma~\ref{l:exposure}.
Lemma~\ref{l:exposure}, followed by Lemma~\ref{l:response-typed},
then gives
\begin{equation}\label{e:rec-exposed}
 a_D(w)|R_D^\tau(w)-R_D^0(w)|
 \le a_D(w)\sum_{i\in\Phi_U(w)}
       \left(\sum_{p\in\mathcal Q_i}\gamma_{i,p}|r_{i,p}|\right)
       \prod_{\substack{j\in\Phi_U(w)\\j\preceq_{\sigma_w}i}}
                           x_{\tau_j}.
\end{equation}
Here \(\mathcal Q_i\) is the set of nonremoved pairs among
the non-arrival neighbors of the departure site \(y=v_i\),
and \(r_{i,p}\) is its pinned-minus-free scalar-product
expectation in the network with \(y\) removed.  With three
present neighbors, \(\gamma_{i,\{f,g\}}\le a_{\tau_i}\) and
\(\gamma_{i,\{f,o\}},\gamma_{i,\{g,o\}}\le b_{\tau_i}\).
With two present neighbors there is one pair and its coefficient
is at most \(b'_{\tau_i}\).  With one present neighbor the
inner sum is zero.  Departures with status \(\mathrm D\) have
ratio one by Definition~\ref{d:ratios} and zero logarithmic
response, so contribute no term.

By Lemma~\ref{l:closure}, removing \(y\) gives an ordinary
domain \(E'_i\), and each \(p\in\mathcal Q_i\) belongs to
\(P_Y(E'_i)\) for its type \(Y\).  For each \(z\in p\)
choose the nearer endpoint of the original pair.  The triangle
inequality gives
\[
 d(z,S)\ge d_0-\min\{d(z,u),d(z,v)\},\qquad
 \min_{z\in p}d(z,S)\ge d_0-s(p)\ge d-s(p).
\]
If \(d-s(p)<0\), use the constant extension of \(F_Y\).
In all cases the definition of the ordinary-domain supremum yields
\begin{equation}\label{e:rec-pair-envelope}
 |r_{i,p}|\le F_Y(d-s(p)).
\end{equation}

We verify every part of the comparison with the lattice kernel.
The removed site \(c\) cannot be visited, and \(w\) first hits
\(v\) at its endpoint, so the lattice symmetry maps it into
\(\mathcal T_n^X\).  A true \(\mathrm U\) departure is a
lattice first visit, since a previous departure would already have
changed its law to \(\mathrm D\) or removed it.  Every lattice
revisit is a \(\mathrm D\) departure.  Some additional lattice
first visits may also be \(\mathrm D\) departures because the
site was already \(\mathrm D\) in the initial domain.  Thus
\[
 N_{\mathrm D}(w)\ge\operatorname{rev}(w),\qquad
 a_D(w)=3^{-n}(3/5)^{N_{\mathrm D}(w)}
             \le3^{-n}(3/5)^{\operatorname{rev}(w)}.
\]
The coefficient identity here is Definition~\ref{d:trace} with
\(N_{\mathrm U}(w)+N_{\mathrm D}(w)=n\); no partition-ratio
factor has been included in \(a_D(w)\).
Passing from the restricted ordering in \eqref{e:rec-exposed}
to the full ordering only inserts factors \(x_{\tau_j}\ge1\)
before a retained index.  It can also insert additional nonnegative
response terms.

A kernel-dead site is removed in the actual trace: this holds for
\(c\), and two departures change any initially nonremoved site
to \(\mathrm X\) at latest.  Consequently every actual response
pair is kernel-live.  At an actual \(\mathrm U\) departure, all
nonremoved non-arrival neighbors have present bonds by the
consistency assertion of Lemma~\ref{l:closure}.  The following
cases exhaust the coefficient comparison.  If both non-outgoing
neighbors are present, they are both kernel-live, and the kernel
assigns exactly the majorants \(a_{\tau_i},b_{\tau_i},b_{\tau_i}\).
If just one is present, its pair with \(o\) has coefficient
at most \(b'_{\tau_i}\).  The kernel assigns that same majorant
when it knows the other neighbor is dead, and assigns
\(b_{\tau_i}\ge b'_{\tau_i}\) if it still regards the other
neighbor as live.  If neither is present there is no response
term.  The outgoing neighbor is always present.  This also
explains why an additional removal cannot increase the required
kernel coefficient.

Using these comparisons in \eqref{e:rec-exposed} and
\eqref{e:rec-pair-envelope}, summing over \(\mathcal W_v\),
enlarging to all kernel trails, and using \(\eta_v\le1\), gives
\begin{equation}\label{e:rec-T1}
 |T_1|\le\sum_{Y=B,C}\sum_{s\ge0}M_{XY,s}F_Y(d-s).
\end{equation}
Spatial lattice symmetries preserve all distances, straight/turn
types, and collinear/diagonal pair types; relabeling sites and pins
preserves the finite product functionals.  Thus the canonical
kernel applies uniformly to the original pair.  Its unshifted
row sums are finite by the (H2)-only summability estimate in the
proof of Lemma~\ref{l:Mtheta}; (H3\(_h\)) is not needed here.

For \(T_2\), Lemma~\ref{l:terminalcover} makes the second
expansion in \(D[w]\), from \(v\) with target \(e\) and
stopping set \(S=S\cup\{e\}\), admissible.  Thus its stopping
set is exactly the one prescribed by Proposition~\ref{p:pairwalk},
and the start is distinct from its nonremoved target.
That proposition
and \(|n_z\cdot n_e|\le1\) give, for either choice of pins,
\[
 |X^\bullet(w)|
 \le\sum_{w'\in\mathcal T_{D[w]}(v,S)}
                a_{D[w]}(w')R_{D[w]}^\bullet(w')
 \le K_*q^{d(v,S)},\qquad
 K_*=(12/5)^{12}\frac{C_V}{1-q}.
\]
The last inequality is the common majorant of
Lemma~\ref{l:onepoint}, so is uniform in the first trail \(w\).
Its hypotheses hold: (H1) and (H2) are assumed, the incoming
bond \(\{c,v\}\) is absent, the start has the required status
and degree, and there are at most twelve exceptional
\(\mathrm U\) departures by Lemma~\ref{l:terminalcover}.
The bound uses the original \(S\), including when the terminal
auxiliary network has extra removed sites.  Both sets of original
legs remain unchanged.  The degree restriction at target \(e\)
was checked in Lemma~\ref{l:terminalcover}; all inserted
observables in this second expansion therefore have separate
degree at most four.

The ratios are positive.  Hence, pointwise in the first trail,
\[
 |R_D^\tau(w)X^\tau(w)-R_D^0(w)X^0(w)|
    \le2\max\{R_D^\tau(w),R_D^0(w)\}K_*q^{d(v,S)}.
\]
The first-trail sum can be enlarged from \(\mathcal W_S\)
to \(\mathcal T_D(u,S)\), dropping its avoidance of \(v\).
The closure lemma applies to that larger class too.  Applying
Lemma~\ref{l:onepoint} is legitimate because its start has absent
incoming bond \(\{c,u\}\), and all its trails have zero
exceptional \(\mathrm U\) departures by Lemma~\ref{l:closure}.
In particular its pointwise assertion for this larger class is
\[
 a_D(w)\max\{R_D^\tau(w),R_D^0(w)\}
          \le(12/5)^{12}\mathsf W(w).
\]
Taking the maximum therefore precedes summing.  Its summed
assertion gives
\[
 |T_2|\le2K_*q^{d(v,S)}
       \sum_{w\in\mathcal T_D(u,S)}
           a_D(w)\max\{R_D^\tau(w),R_D^0(w)\}
 \le2K_*^2q^{d(u,S)+d(v,S)}\le C_Tq^{2d}.
\]
Combining this with \eqref{e:rec-T1} and taking the supremum
over the defining class of \(F_X(d)\) proves the result.  If
that class is empty its supremum is zero and the same inequality
holds by nonnegativity of the right side.
\end{proof}

In the next lemma \(X,Y\) range over \(\{B,C\}\), and every
shift index \(s\) ranges over the nonnegative integers.

\begin{lemma}\label{l:renewal}
Let \(m_{XY,s}\ge0\), \(h_X>0\), \(0<q^2<\theta<1\),
\(0<\rho<1\), and
\(\sum_{Y,s}m_{XY,s}h_Y\theta^{-s}\le\rho h_X\).
Suppose \(0\le F_X(d)\le2\) on \(\mathbb Z\), with
\(F_X(d)=F_X(0)\) for \(d<0\), and, for integers \(d\ge1\),
\[
 F_X(d)\le\sum_{Y,s}m_{XY,s}F_Y(d-s)+C_Tq^{2d},
 \qquad C_T\ge0.
\]
Writing \(H=\max_Xh_X/\min_Xh_X\), one has, for \(d\ge0\),
\[
 F_X(d)\le A\theta^d,\qquad
 A=\frac{H\max\{2,C_T\}}{1-\rho}\ge2 .
\]
\end{lemma}

\begin{proof}
For a nonnegative function \(f\) on \(\{B,C\}\times\mathbb Z\),
define the nonnegative shift operator
\[
 (Kf)_X(d)=\sum_{Y,s}m_{XY,s}f_Y(d-s).
\]
For an arbitrary nonnegative \(f\), this denotes the extended
nonnegative sum; every application used below is finite by the
unshifted or shifted row bound, as specified in the proof.
Put \(h_{\min}=\min_Xh_X>0\), \(C_*=\max\{2,C_T\}\), and
\[
 g_X(d)=
 \begin{cases}2,&d\le0,\\ C_Tq^{2d},&d\ge1.\end{cases}
\]
The assumed inequality and \(F_X(d)\le2\) give \(F\le KF+g\)
on all integers.  Also
\(\sum_{Y,s}m_{XY,s}h_Y\le\rho h_X\), since
\(\theta^{-s}\ge1\).  In particular all row sums of \(K\)
are finite, so it acts on bounded functions and preserves
nonnegativity and inequalities.  Iterating a finite number of
times yields, for every integer \(k\ge1\),
\begin{equation}\label{e:rec-renewal-iterate}
 F\le K^kF+\sum_{j=0}^{k-1}K^jg.
\end{equation}

View \(h_X\) as constant in \(d\).  The unshifted row inequality
and induction show \(K^kh\le\rho^kh\).  Since
\(F_X(d)\le2h_X/h_{\min}\) for every integer \(d\),
\begin{equation}\label{e:rec-renewal-remainder}
 (K^kF)_X(d)\le\frac{2}{h_{\min}}\rho^kh_X.
\end{equation}
This is an unweighted bound, valid even when the iterated shifts
reach arbitrarily negative integers.

For the source, if \(d\le0\) then \(\theta^d\ge1\), and if
\(d\ge1\) then \(q^{2d}\le\theta^d\).  Thus for every integer
\(d\),
\[
 g_X(d)\le\frac{C_*}{h_{\min}}h_X\theta^d.
\]
Writing \(e_X(d)=h_X\theta^d\), the shifted row inequality is
exactly \(Ke\le\rho e\).  Its sums are finite even for
negative \(d\), by the stated weighted summability hypothesis.
Induction and positivity therefore imply
\[
 (K^jg)_X(d)\le\frac{C_*}{h_{\min}}\rho^jh_X\theta^d.
\]
All rearrangements used here involve nonnegative summands; they
can equivalently be proved first for finite shift sums and then
by increasing limits.  Combining these bounds with
\eqref{e:rec-renewal-iterate}, for fixed \(d\), gives
\[
 F_X(d)\le\frac{2h_X}{h_{\min}}\rho^k
    +\frac{C_*h_X}{h_{\min}}\theta^d
                       \frac{1-\rho^k}{1-\rho}.
\]
Let \(k\) tend to infinity.  The first term tends to zero
because \(0<\rho<1\).  Since \(h_X/h_{\min}\le H\),
the desired bound follows for every \(d\ge0\).
Finally \(H\ge1\), \(C_*\ge2\), and \(0<1-\rho<1\),
so the displayed \(A\) is at least two.
\end{proof}

\begin{theorem}\label{t:rec}
Fix \((c_{\rm s},c_{\rm t})=(961/1000,187/200)\).
Assume (H1), (H2) with any \(C_V\ge1\), \(0<q<1\), and
(H3\(_h\)) for the type-dependent kernel and fixed orderings above.
Then there exist \(A\ge2\) and \(0<\theta<1\) such that
\(\mathsf P(A,\theta)\) holds.
One may take the \(\theta,\rho,h\) supplied by the shifted-witness
lemma and
\[
 A=\frac{\max h}{\min h}\,
       \frac{\max\{2,\,2(12/5)^{24}(C_V/(1-q))^2\}}{1-\rho}.
\]
\end{theorem}

\begin{proof}
Lemma~\ref{l:Mtheta} supplies \(q^2<\theta<1\) and
\(\rho_0<\rho<1\), with the same fixed positive vector \(h\),
such that \(M_\theta h\le\rho h\).  Set
\(m_{XY,s}=M_{XY,s}\).  The envelope bounds and constant
extension are part of Definition~\ref{d:pairenv};
Lemma~\ref{l:recineq} supplies the recursion on all positive
integers with finite
\(C_T=2(12/5)^{24}(C_V/(1-q))^2\).
Lemma~\ref{l:renewal} therefore gives
\(F_X(d)\le A\theta^d\) for both pair types and all integers
\(d\ge0\), with exactly the stated finite \(A\ge2\).

To obtain the property in Definition~\ref{h:P}, fix its data
\(G,\tau,Y,u,v\).  Set all sites of \(Y\) to \(\mathrm X\)
and all other sites to \(\mathrm U\) with their full kept sets.
This is an ordinary consistent domain: \(Y\cap S=\varnothing\),
and every bond between nonremoved sites is kept by both ends.
It has precisely the partition sums of
Definition~\ref{d:pinmodel}, with the original legs and no new
pins at removed sites.  Designate the given common neighbor in
\(Y\) as center.  A distance-two lattice pair is either collinear
or diagonal, so it lies in \(P_B\) or \(P_C\).
The bound for the corresponding envelope, at
\(d=\min\{d(u,S),d(v,S)\}\), is the required inequality.
If there are no such pairs, the assertion is vacuous.  In the
nonvacuous case \(G\ne\varnothing\) is finite.  A site with
maximal first coordinate has a neighbor outside \(G\), so
\(S\ne\varnothing\) and this minimum is a nonnegative integer.
Thus all quantifiers
of \(\mathsf P(A,\theta)\) are satisfied.
\end{proof}
